\chapter{DFS temporal guiado por score}
\label{sec:paper-dfs}

\section{Rol dentro del pipeline}

El matcher SIMD (cap.~\ref{sec:paper-simd}) enumera paths
\emph{estructuralmente} válidos: cumplen los tipos de nodo/arista
y la longitud del patrón. Sobre ese conjunto se aplica una segunda
fase, \emph{refinement}, que:

\begin{enumerate}
    \item impone orden temporal sobre las aristas del path;
    \item computa score compuesto por path;
    \item poda paths debajo de umbral o que exceden cotas de
    ramificación (\code{k-simplicity}).
\end{enumerate}

La fase \emph{refinement} se implementa como un DFS temporal
guiado por score parcial. Este capítulo describe el algoritmo, las
invariantes que preserva, y por qué usamos DFS en vez de las
variantes clásicas (Ullmann, VF2, TurboISO).

\section{Hipótesis como forma de patrón}

Una hipótesis en el DSL tiene forma:

\begin{lstlisting}[basicstyle=\ttfamily\small]
User -[Auth]-> Host -[Connect]-> IP {k=3, window=30m, min_score=0.4}
\end{lstlisting}

que se desazucara a una secuencia ordenada de restricciones:

\begin{itemize}
    \item tipos por posición $(\tau_0, \tau_1, \tau_2, \tau_3)$;
    \item tipos de arista por salto $(\rho_0, \rho_1, \rho_2)$;
    \item ventana temporal $W$ sobre el path completo;
    \item cota estructural $k$ (rama del DFS);
    \item filtro post-hoc $\text{score} \ge s_{\min}$.
\end{itemize}

\begin{definition}[Path instancia del patrón]
Una secuencia $\pi = (v_0, e_0, v_1, e_1, \dots, v_n)$ es instancia
si: (i) $\text{type}(v_i) = \tau_i$ para $i = 0..n$; (ii)
$\text{type}(e_i) = \rho_i$ para $i = 0..n-1$; (iii)
$t(e_0) \le t(e_1) \le \dots \le t(e_{n-1})$ y
$t(e_{n-1}) - t(e_0) \le W$; (iv) la ramificación de cada nodo
interno del patrón es $\le k$.
\end{definition}

\section{Algoritmo}

\begin{algorithm}[H]
\caption{DFS temporal guiado por score parcial}
\label{alg:paper-dfs}
\KwIn{Grafo $G$, patrón $(\tau, \rho)$, ventana $W$, $k$, $s_{\min}$}
\KwOut{Paths ordenados por score}
$R \gets \emptyset$ \Comment{resultado}
\ForEach{$v_0 \in V: \text{type}(v_0) = \tau_0$}{
    DFS($v_0$, path$=[v_0]$, $t_0=-\infty$, score$=0$, depth$=0$)\;
}
\Return top $N$ de $R$ por score\;

\BlankLine
\textbf{procedure} DFS($v$, path, $t_{\min}$, score, depth):\\
\If{depth $= n$}{
    \If{$\text{score} \ge s_{\min}$}{
        $R \gets R \cup \{(\text{path}, \text{score})\}$\;
    }
    \Return\;
}
$\tau_{\text{next}} \gets \tau_{\text{depth}+1}$;
$\rho_{\text{next}} \gets \rho_{\text{depth}}$\;
$N_v \gets \{(e, u) : e = (v, u), \text{type}(e) = \rho_{\text{next}}, \text{type}(u) = \tau_{\text{next}}\}$\;
\If{$|N_v| > k$}{\Return \Comment{poda por k-simplicity}}
\ForEach{$(e, u) \in N_v$}{
    \If{$t(e) < t_{\min} \lor t(e) - t(e_0) > W$}{continue}
    $s' \gets \text{update\_partial\_score}(\text{score}, v, e, u)$\;
    \If{$\text{upper\_bound}(s', \text{depth}+1) < s_{\min}$}{
        continue \Comment{poda por score}
    }
    DFS($u$, path $\cup \{e, u\}$, $t(e)$, $s'$, depth+1)\;
}
\end{algorithm}

\section{Propiedades y complejidad}

\begin{lemma}[Orden temporal]
Todo path $\pi$ emitido por DFS satisface
$t(e_0) \le t(e_1) \le \dots$ por construcción
(gracias a $t_{\min}$).
\end{lemma}

\begin{lemma}[Monotonía del score parcial]
Si $\text{upper\_bound}(s, d)$ acota superiormente el score
alcanzable extendiendo $d$ hops más, podar cuando
$\text{upper\_bound} < s_{\min}$ es seguro: ningún path emitido por
esa rama superaría $s_{\min}$.
\end{lemma}

La cota superior se construye asumiendo puntajes máximos locales
para los componentes que monotonizan (entity rarity, temporal
novelty) y cero para los que requieren contexto completo
(neighborhood concentration). No es apretada pero sí válida.

La complejidad en el peor caso es
$O(|V_{\tau_0}| \cdot k^n)$. En la práctica el DFS emite orden de
decenas a miles de paths para hipótesis de 3-4 hops sobre grafos
de $10^4$-$10^5$ nodos, con tiempos sub-segundo reportados en
benchmarks.

\section{¿Por qué DFS y no Ullmann/VF2?}

Los clásicos (Ullmann~\cite{ullmann1976}, VF2~\cite{cordella2004vf2},
TurboISO) resuelven \emph{subgraph isomorphism} sobre grafos
estáticos con pesos uniformes. Nuestro setup difiere en tres
puntos:

\begin{enumerate}
    \item \textbf{Multigrafo temporal}. Múltiples aristas entre un
    par de nodos distinguidas por timestamp. Ullmann/VF2 asumen
    único edge.
    \item \textbf{Score como guía}, no como post-filtro. Queremos
    podar ramas prometedoramente malas temprano. DFS con
    \emph{best-first} en la frontera de siblings es natural.
    \item \textbf{Paths, no subgrafos}. Buscamos caminos lineales
    (a veces con fork), no cliques ni cualquier subgrafo.
\end{enumerate}

DFS encaja los tres. La complejidad es peor en el peor caso pero
la cota $k$ (\emph{k-simplicity}) lo acota estructuralmente y el
upper-bound de score lo acota semánticamente.
